% ------------------------------------------------------------------------------
% Introdução
% ------------------------------------------------------------------------------

\chapter{Introdução}
\label{chap_introducao}



A queda livre é um dos fenômenos mais clássicos e fundamentais da física, tendo desempenhado papel central no desenvolvimento da ciência moderna. Desde Galileu Galilei, que desafiou concepções aristotélicas ao demonstrar experimentalmente que corpos de diferentes massas caem com a mesma aceleração na ausência de resistência do ar, o estudo desse movimento impulsionou avanços teóricos e experimentais na mecânica.

O interesse pela queda livre vai além do contexto histórico: trata-se de um problema paradigmático para a modelagem matemática, a análise de métodos numéricos e a validação de soluções analíticas. Em aplicações modernas, compreender a influência da resistência do ar é essencial em áreas como engenharia aeroespacial, balística, meteorologia, esportes e ensino de ciências.

Apesar de sua aparente simplicidade, a queda livre apresenta desafios relevantes quando fatores como resistência do ar linear e quadrática são considerados. A análise desses efeitos permite não apenas aprimorar a compreensão dos fenômenos físicos, mas também desenvolver e validar algoritmos computacionais robustos para a solução de equações diferenciais ordinárias.

Este trabalho tem como objetivo investigar, por meio de simulações computacionais, o comportamento de corpos em queda livre sob diferentes regimes físicos: sem resistência do ar, com resistência linear (Lei de Stokes) e com resistência quadrática. São utilizados métodos numéricos e soluções analíticas para avaliar a precisão dos algoritmos implementados, comparar os resultados obtidos e discutir as limitações e potencialidades de cada abordagem.

Além de contribuir para a compreensão dos aspectos teóricos e práticos do problema, este estudo busca demonstrar a importância da modelagem matemática e da análise computacional na solução de problemas reais, destacando a influência dos parâmetros físicos e numéricos na acurácia dos resultados. Espera-se, ao final, fornecer uma visão abrangente e aplicada sobre o tema, evidenciando sua relevância para a formação científica, tecnológica e didática.

